\documentclass[10pt,aspectratio=169]{beamer}

\usetheme{Madrid}
\usecolortheme{whale}
\setbeamertemplate{navigation symbols}{}
\setbeamertemplate{caption}[numbered]

\usepackage[utf8]{inputenc}
\usepackage[T1]{fontenc}
\usepackage[english]{babel}
\usepackage{amsmath,amssymb}
\usepackage{graphicx}
\usepackage{booktabs}

\title[Cafeteria simulation]{Performance Modeling of a Cafeteria System}
\subtitle{A discrete-event simulation study}
\author{Daniel}
\institute[Tor Vergata]{Performance Modeling of Computer Systems and Networks \\
                        Universit\`a degli Studi di Roma Tor Vergata}
\date{A.A.\ 2025/2026}

\begin{document}

% ====================================================================
\begin{frame}
\titlepage
\end{frame}

% ====================================================================
\begin{frame}{Outline}
\tableofcontents
\end{frame}

% ====================================================================
\section{Motivation and system}
% ====================================================================
\begin{frame}{Motivation}
\textbf{Operational dilemma of a small cafeteria:}
\begin{itemize}
  \item Demand grows -- should the manager hire a \alert{waiter} or buy a \alert{seat}?
  \item Intuition: ``add a waiter when the queue grows''.
  \item Intuition is \alert{not necessarily correct}: a customer occupies a seat throughout the entire stay, including consumption.
\end{itemize}
\vspace{0.5cm}
\textbf{Three concrete questions answered in this study:}
\begin{enumerate}
  \item Which is the system's bottleneck, the waiters or the seats?
  \item How sensitive are performance metrics to the arrival rate?
  \item Given operating costs and a revenue model, what is the operational policy that maximises the hourly net benefit per demand level?
\end{enumerate}
\end{frame}

% ====================================================================
\begin{frame}{System description}
\begin{columns}[T]
\begin{column}{0.5\textwidth}
\textbf{Two customer classes (Bernoulli):}
\begin{itemize}
  \item \texttt{CAFE}: $p_c = 0.7$
  \item \texttt{DESAYUNO}: $p_d = 0.3$
\end{itemize}

\vspace{0.4cm}
\textbf{Two resources:}
\begin{itemize}
  \item $S$ \textbf{seats} (held for the whole stay)
  \item $C$ \textbf{waiters} (held only during attention)
\end{itemize}

\vspace{0.4cm}
\textbf{Two FCFS queues:}
\begin{itemize}
  \item Standing queue (unbounded)
  \item Seated-waiting queue (size $\le S$)
\end{itemize}
\end{column}
\begin{column}{0.5\textwidth}
\textbf{Customer lifecycle:}

\begin{small}
\fbox{\parbox{\linewidth}{
\textsc{Arrival} \\
$\downarrow$ (if no seat) \\
standing queue \\
$\downarrow$ seat allocated \\
seated-waiting \\
$\downarrow$ waiter assigned \\
\textsc{Attention begins} \\
$\downarrow$ \textsc{End of attention} \\
consuming \\
$\downarrow$ \textsc{End of consumption} \\
\textsc{Departure} (seat freed)
}}
\end{small}
\end{column}
\end{columns}
\end{frame}

% ====================================================================
\begin{frame}{Why simulation?}
\textbf{In isolation:} each resource would be an M/M/$c$ system.

\vspace{0.3cm}
\textbf{In reality:} the seat is also held \alert{while the customer waits for a waiter}.
\begin{itemize}
  \item The two resources are \alert{coupled}.
  \item No closed-form expression exists for the joint distribution.
\end{itemize}

\vspace{0.5cm}
\textbf{Methods used (all from course material):}
\begin{itemize}
  \item Algorithm 1.1 of Leemis \& Park for the model-development workflow
  \item Algorithm 1.2 of Leemis \& Park for the next-event simulation engine
  \item Multi-stream Lehmer generator for pseudo-random numbers
  \item Welch's method for warm-up identification
  \item Replication-deletion with t-Student CI for the statistical analysis
\end{itemize}
\end{frame}

% ====================================================================
\section{Specification model}
% ====================================================================
\begin{frame}{Specification model: parameters}
\begin{table}
\centering
\begin{tabular}{l l c}
\toprule
Process & Distribution & Mean (min) \\
\midrule
Inter-arrivals          & Exponential & $1/\lambda$ (varied) \\
Class assignment        & Bernoulli   & $p_c = 0.7$ (coffee) \\
Attention -- coffee     & Exponential & $E[A_c] = 1.5$ \\
Attention -- breakfast  & Exponential & $E[A_d] = 4$ \\
Consumption -- coffee   & Exponential & $E[C_c] = 5$ \\
Consumption -- breakfast& Exponential & $E[C_d] = 12$ \\
\bottomrule
\end{tabular}
\end{table}

\vspace{0.3cm}
\textbf{Why exponential service times?}
\begin{itemize}
  \item Memoryless property $\Rightarrow$ preserves Markov property in sub-models
  \item Allows analytical bounds for validation
  \item Standard reference in queueing theory
\end{itemize}
\end{frame}

% ====================================================================
\begin{frame}{Derived means and analytical bounds}
\textbf{Marginal means averaged over the two classes:}
\begin{align*}
E[A] &= p_c\,E[A_c] + p_d\,E[A_d] = 0.7\cdot 1.5 + 0.3\cdot 4 = \alert{2.25 \text{ min}} \\
E[C] &= p_c\,E[C_c] + p_d\,E[C_d] = 0.7\cdot 5 + 0.3\cdot 12 = \alert{6.5 \text{ min}} \\
E[T_{\text{seat}}] &= E[A] + E[C] = \alert{8.75 \text{ min}}
\end{align*}

\vspace{0.3cm}
\textbf{Analytical utilization bounds} (treating each resource in isolation):
\[
\hat\rho_{\text{seat}}(C,S,\lambda) = \frac{\lambda\,E[T_{\text{seat}}]}{S}, \qquad
\hat\rho_{\text{waiter}}(C,S,\lambda) = \frac{\lambda\,E[A]}{C}
\]

\vspace{0.3cm}
\textbf{Key observation:} $E[T_{\text{seat}}] \gg E[A]$ $\Rightarrow$ the seat saturates first.

\hfill\alert{The seat is the predicted bottleneck.}
\end{frame}

% ====================================================================
\section{Simulator and verification}
% ====================================================================
\begin{frame}{Computational model}
\textbf{Implemented in ANSI C, following Algorithm 1.2 of Leemis \& Park.}

\vspace{0.4cm}
\textbf{Next-event main loop:}
\begin{enumerate}
  \item Initialize, schedule first arrival
  \item Find next event = $\min\{t_{\text{arrival}}, t_{\text{end-att}}[i], t_{\text{end-cons}}[j]\}$
  \item Advance clock to that time; update statistics; process event
  \item Repeat until $\text{clients\_served} \ge 10\,000$
\end{enumerate}

\vspace{0.3cm}
\textbf{Multi-stream Lehmer RNG:}
\begin{itemize}
  \item $m = 2^{31}-1$, $a = 48271$, 256 disjoint streams, jump multiplier $22\,925$
  \item \alert{One stream per stochastic process} (decoupling):
  \begin{itemize}
    \item Stream 0: inter-arrivals \quad
          Stream 1: class
    \item Streams 2-3: attention (coffee/breakfast)
    \item Streams 4-5: consumption (coffee/breakfast)
  \end{itemize}
  \item Variate generation by inverse method: $X = -m\,\ln(1 - U)$
\end{itemize}
\end{frame}

% ====================================================================
\begin{frame}{Runtime verification}
\textbf{Three distribution-free checks confirm the simulator:}

\vspace{0.5cm}
\textbf{(i) Little's Law:} $L = \lambda\,W$ holds in every replication.
\[
\text{Baseline: } \lambda\,W = \tfrac{1}{3} \cdot 16.06 = 5.36 = L \quad \checkmark
\]

\vspace{0.3cm}
\textbf{(ii) Observed arrival rate} $= 1 / E[\text{IAR}]$.
\[
\text{Baseline: observed } X = 20.01 \text{ cust/h} = \lambda \quad \checkmark
\]

\vspace{0.3cm}
\textbf{(iii) Class mix} converges to the input probabilities.
\[
69.9\% \text{ coffee} / 30.1\% \text{ breakfast} \approx 70/30 \quad \checkmark
\]

\vspace{0.3cm}
\hfill\textit{Three independent checks $\Rightarrow$ the simulator is verified.}
\end{frame}

% ====================================================================
\section{Transient analysis (Welch)}
% ====================================================================
\begin{frame}{Welch's method -- warm-up identification}
\textbf{Problem:} the simulator starts empty $\Rightarrow$ early observations are biased.

\vspace{0.4cm}
\textbf{Procedure} (Leemis \& Park, lect.\ 14):
\begin{enumerate}
  \item Run $R = 30$ independent replications; take snapshots every $\Delta t = 60$ min.
  \item Ensemble average across replications at each instant $k$:
  \[
  \bar Y_k = \frac{1}{R}\sum_{r=1}^{R} Y_{r,k}
  \]
  \item Smooth with a moving-average filter of window $2w+1$ ($w = 5 \Rightarrow$ 11 points):
  \[
  \bar Y_k^{(w)} = \frac{1}{2w+1}\sum_{s=-w}^{w} \bar Y_{k+s}
  \]
  \item Visually identify $T_0$ = the instant when the smoothed curve stabilises.
\end{enumerate}

\vspace{0.3cm}
\textbf{Applied to 3 different indicators:} cumulative $W$, instantaneous $n$ in system,
time-integrated seat utilization.
\end{frame}

% ====================================================================
\begin{frame}{Welch's method -- result}
\begin{columns}[T]
\begin{column}{0.5\textwidth}
\centering
\includegraphics[width=\linewidth]{welch_avg_total.png}
\\[0.2cm]
{\small Cumulative mean response time}
\end{column}
\begin{column}{0.5\textwidth}
\centering
\includegraphics[width=\linewidth]{welch_util_chairs.png}
\\[0.2cm]
{\small Time-integrated seat utilization}
\end{column}
\end{columns}

\vspace{0.4cm}
\begin{center}
\textbf{All three indicators stabilise by} $t \approx 50$ h.
\\[0.2cm]
\fbox{\quad $T_0 = 60$ h with a 10\,h safety margin\quad}
\\[0.2cm]
All subsequent statistics computed over $[T_0, T_{\text{end}}]$.
\end{center}
\end{frame}

% ====================================================================
\section{Experimental design and statistics}
% ====================================================================
\begin{frame}{Statistical workflow and confidence intervals}
\textbf{Two-stage study:}
\begin{enumerate}
\item \textbf{Warm-up identification} (Welch, slide 11): \\
      $R = 30$ short replications on baseline $\Rightarrow T_0 = 60$ h
\item \textbf{Steady-state estimation} (replication-deletion, deleting $[0,T_0]$):
   \begin{itemize}
     \item \textbf{Baseline} $(C,S,\lambda)=(2,4,20)$ in detail: \alert{$n = 64$} replications
     \item \textbf{Full factorial} $(C,S)\in\{2,3\}\times\{4,5\}$, $\lambda\in\{15,20,25,30\}$:\\
           \alert{$n = 32$} replications $\times$ 16 scenarios $=$ 512 runs
   \end{itemize}
\end{enumerate}

\vspace{0.3cm}
\textbf{95\% t-Student CI:}
\[
\bar X \;\pm\; t_{n-1,\,0.025}\,\frac{s}{\sqrt n}, \qquad
t_{63} \approx 1.998 \text{ (baseline)}, \quad
t_{31} \approx 2.040 \text{ (factorial)}
\]

\vspace{0.3cm}
\textbf{Why replication-deletion?} Replications are \alert{strictly independent} by construction
(disjoint Lehmer streams) — no auto-correlation analysis needed.
\end{frame}

% ====================================================================
\section{Results}
% ====================================================================
\begin{frame}{Baseline result: $(C,S,\lambda) = (2,4,20)$, $n = 64$ replications}
\begin{table}
\centering
\begin{tabular}{l r r r}
\toprule
Metric & Mean & St.\ dev. & 95\% CI half-width \\
\midrule
$W$ (min)                       & 16.06 & 1.213 & $\pm$ 0.30 (1.9\%) \\
Standing wait (min)             &  6.54 & 1.149 & $\pm$ 0.29 (4.4\%) \\
Seated wait (min)               &  0.15 & 0.010 & $\pm$ 0.002 (1.6\%) \\
Throughput (cust/h)             & 20.01 & 0.236 & $\pm$ 0.06 (0.3\%) \\
Seat utilization                & 0.7935 & 0.014 & $\pm$ 0.0035 (0.4\%) \\
Waiter utilization              & 0.3765 & 0.007 & $\pm$ 0.0017 (0.5\%) \\
Mean number in system $L$       &  5.36 & 0.444 & $\pm$ 0.11 (2.1\%) \\
\bottomrule
\end{tabular}
\end{table}

\vspace{0.4cm}
\textbf{Sanity checks:}
\begin{itemize}
  \item Little: $\lambda W = \tfrac{1}{3}\cdot 16.06 = 5.36 = L$ \quad\checkmark
  \item $X = 20.01 \approx \lambda$ \quad (no losses, unbounded queue) \checkmark
  \item All half-widths $< 4.5\%$ $\Rightarrow$ excellent statistical precision.
\end{itemize}
\end{frame}

% ====================================================================
\begin{frame}{Response time across all scenarios}
\begin{columns}[T]
\begin{column}{0.5\textwidth}
\centering
\includegraphics[width=\linewidth]{fig_W_vs_lambda.png}
\\[0.2cm]
{\small linear scale}
\end{column}
\begin{column}{0.5\textwidth}
\centering
\includegraphics[width=\linewidth]{fig_W_vs_lambda_log.png}
\\[0.2cm]
{\small logarithmic scale}
\end{column}
\end{columns}

\vspace{0.3cm}
\begin{itemize}
\item Stable configurations: $W \in [10, 35]$ min across the whole range.
\item $(2,4)$ and $(3,4)$ \alert{saturate} at $\lambda = 30$/h $\Rightarrow$ $W$ explodes.
\item The log scale highlights the order-of-magnitude difference.
\end{itemize}
\end{frame}

% ====================================================================
\begin{frame}{Bottleneck identification}
\begin{columns}[T]
\begin{column}{0.5\textwidth}
\centering
\includegraphics[width=\linewidth]{fig_util_chairs.png}
\\[0.1cm]
{\small \textbf{Seat utilization}}
\end{column}
\begin{column}{0.5\textwidth}
\centering
\includegraphics[width=\linewidth]{fig_util_waiters.png}
\\[0.1cm]
{\small Waiter utilization (note same scale)}
\end{column}
\end{columns}

\vspace{0.3cm}
\textbf{The seat is the bottleneck at every operating point:}
\begin{itemize}
\item Baseline ($\lambda = 20$): $\rho_{\text{seat}} \approx 0.79$ vs $\rho_{\text{waiter}} \approx 0.38$
\item Peak ($\lambda = 30$): seat \alert{saturates at 1.0}, waiter stays below 0.56
\item Confirms the analytical prediction $E[T_{\text{seat}}] \gg E[A]$.
\end{itemize}
\end{frame}

% ====================================================================
\begin{frame}{Effect of adding one extra resource (at $\lambda = 25$/h)}
\begin{table}
\centering
\begin{tabular}{l r r}
\toprule
Configuration & $W$ (min) & Reduction vs baseline \\
\midrule
$(C,S) = (2,4)$ baseline   & 162.6 & --- \\
$(C,S) = (2,5)$ +1 seat    & \alert{15.3} & \alert{$-90\%$} \\
$(C,S) = (3,4)$ +1 waiter  & 84.0  & $-48\%$ \\
$(C,S) = (3,5)$ both       & 13.6  & $-92\%$ \\
\bottomrule
\end{tabular}
\end{table}

\vspace{0.6cm}
\begin{center}
\Large
\alert{One extra seat $\to$ 90\% lower $W$.} \\
\alert{One extra waiter $\to$ only 48\% lower.}
\end{center}

\vspace{0.4cm}
The waiter upgrade is largely redundant: the bottleneck is so clearly the seat that
$(3,5)$ barely improves over $(2,5)$ at this rate.
\end{frame}

% ====================================================================
\section{Validation}
% ====================================================================
\begin{frame}{Validation against analytical predictions}
\begin{columns}[T]
\begin{column}{0.5\textwidth}
\textbf{Waiter side:}

$\rho_{\text{waiter}}^{\text{sim}}$ agrees with $\hat\rho_{\text{waiter}}$
within \alert{$\le 0.5$ pp} across all 16 scenarios.

\vspace{0.3cm}
$\Rightarrow$ Validates the waiter sub-model: a waiter behaves exactly as an isolated M/M/$C$.
\end{column}
\begin{column}{0.5\textwidth}
\textbf{Seat side:}

$\rho_{\text{seat}}^{\text{sim}}$ is systematically \alert{5--8 pp higher} than $\hat\rho_{\text{seat}}$.

\vspace{0.3cm}
\textbf{Reason:} the bound ignores the time spent seated waiting for a waiter:
\[
E[T_{\text{seat}}^{\text{real}}] = E[T_{\text{seat}}] + E[W_{\text{seated-waiting}}]
\]
\end{column}
\end{columns}

\vspace{0.4cm}
\begin{center}
\fbox{\parbox{0.85\textwidth}{\centering
\alert{This coupling is precisely what defeats a closed-form solution
and justifies the use of simulation.}
}}
\end{center}

\vspace{0.3cm}
\textbf{Stability prediction:} when $\hat\rho_{\text{seat}} > 1$, simulation
shows saturation ($\rho^{\text{sim}} = 1.000$, $W$ explodes) $\Rightarrow$
the bound is qualitatively correct.
\end{frame}

% ====================================================================
\section{Economic analysis}
% ====================================================================
\begin{frame}{Hourly net benefit}
\textbf{Benefit function} (cost / revenue per hour):
\[
B(C,S,\lambda) = \underbrace{\big(p_c\,p_{\text{coffee}} + p_d\,p_{\text{breakfast}}\big)\,X(C,S,\lambda)}_{\text{revenue}}
                 \;-\; \underbrace{c_w\,C + c_s\,S}_{\text{operating cost}}
\]

\textbf{Parameters:}
\begin{itemize}
  \item Prices: $p_{\text{coffee}} = 1.5$ \euro, $p_{\text{breakfast}} = 5$ \euro
        $\Rightarrow$ $\bar r = 2.55$ \euro/customer
  \item Costs: $c_w = 12$ \euro/h, $c_s = 1$ \euro/h
\end{itemize}

\vspace{0.4cm}
\begin{center}
\Large \alert{A seat is 12$\times$ cheaper than a waiter.}
\end{center}

\vspace{0.3cm}
The chair cost amortises the capital expense (purchase, floor space, maintenance)
over operating hours, allowing a like-for-like comparison with the waiter salary.
\end{frame}

% ====================================================================
\begin{frame}{Operational policy}
\begin{columns}[T]
\begin{column}{0.55\textwidth}
\centering
\includegraphics[width=\linewidth]{fig_benefit_vs_lambda.png}
\end{column}
\begin{column}{0.45\textwidth}
\textbf{Recommended policy as a function of demand:}
\begin{itemize}
  \item $\lambda = 15$/h $\to$ \alert{$(2,4)$}: cheapest, fine performance.
  \item $\lambda = 20$/h $\to$ \alert{$(2,4)$}: baseline still wins.
  \item $\lambda = 25$/h $\to$ \alert{$(2,5)$}: \textit{one extra seat is enough}.
  \item $\lambda = 30$/h $\to$ \alert{$(3,5)$}: both required, even so $W = 33.7$ min.
\end{itemize}

\vspace{0.4cm}
\begin{center}
\alert{\textbf{Rule:} add a seat before adding a waiter.}
\end{center}
\end{column}
\end{columns}
\end{frame}

% ====================================================================
\section{Conclusions}
% ====================================================================
\begin{frame}{Conclusions}
\textbf{Main findings:}
\begin{enumerate}
  \item \alert{The seat -- not the waiter -- is the bottleneck} of the system at
        every operating point. The analytical bound $\hat\rho_{\text{seat}} = \lambda E[T_{\text{seat}}]/S$
        correctly predicts both the magnitude and the stability boundary.
  \item Adding a seat is operationally and economically far superior to adding a
        waiter: \alert{12$\times$ cheaper} and directly addresses the bottleneck.
  \item The optimal configuration is \alert{demand-dependent}: $(2,4)$ at low to
        normal demand, $(2,5)$ at high demand, $(3,5)$ only at peak.
\end{enumerate}

\vspace{0.4cm}
\textbf{Limitations and possible extensions:}
\begin{itemize}
  \item Service times assumed exponential (standard but not always accurate).
  \item Customer balking and reneging not modelled.
  \item Extensions: heterogeneous waiters, time-varying $\lambda$ (peak hours),
        take-away service not needing a seat.
\end{itemize}
\end{frame}

% ====================================================================
\begin{frame}
\begin{center}
\Huge Thank you. \\
\vspace{1cm}
\Large Questions?
\end{center}
\end{frame}

\end{document}
